% ======================================================================
% CONFIGURACIÓN Y PREÁMBULO GLOBAL - MEMORIA TFG ETSIIT - UGR
% ======================================================================

% ---- Codificación e Idioma ----
\usepackage[utf8]{inputenc}
\usepackage[T1]{fontenc}
\usepackage[spanish,es-tabla,es-nodecimaldot]{babel}

% ---- Tipografía y Márgenes ----
\usepackage{lmodern}
\usepackage[a4paper, top=2.5cm, bottom=2.5cm, left=3.0cm, right=2.5cm, headheight=16pt]{geometry}
\renewcommand{\baselinestretch}{1.18} % Interlineado académico profesional

% ---- Matemáticas y Símbolos ----
\usepackage{amsmath, amssymb}
\usepackage{xspace}
\providecommand{\bm}[1]{\mathbf{#1}}

% ---- Tablas y Figuras ----
\usepackage{graphicx}
\usepackage{booktabs}
\usepackage{multirow}
\usepackage{longtable}
\usepackage{tabularx}
\usepackage{array}
\usepackage{float}
\usepackage{caption}
\usepackage{subcaption}

% ---- Encabezados y Pies de Página ----
\usepackage{fancyhdr}
\pagestyle{fancy}
\fancyhf{}
\fancyhead[LE,RO]{\thepage}
\fancyhead[RE]{\small\nouppercase{\leftmark}}
\fancyhead[LO]{\small\nouppercase{\rightmark}}
\renewcommand{\headrulewidth}{0.4pt}
\renewcommand{\footrulewidth}{0pt}

% ---- Colores Corporativos y de Diseño ----
\usepackage{xcolor}
\definecolor{ugrred}{RGB}{198, 12, 48}       % Rojo institucional UGR
\definecolor{etsiitblue}{RGB}{0, 68, 124}     % Azul institucional ETSIIT
\definecolor{darkblue}{RGB}{12, 45, 110}
\definecolor{darkgreen}{RGB}{20, 110, 45}
\definecolor{codebg}{RGB}{248, 249, 250}
\definecolor{codeframe}{RGB}{220, 224, 230}
\definecolor{codekeyword}{RGB}{175, 0, 80}
\definecolor{codestring}{RGB}{20, 130, 40}
\definecolor{codecomment}{RGB}{110, 115, 125}

% ---- Formato de Código Fuente ----
\usepackage{listings}
\lstdefinestyle{pythonstyle}{
    language=Python,
    backgroundcolor=\color{codebg},
    basicstyle=\ttfamily\footnotesize,
    keywordstyle=\color{codekeyword}\bfseries,
    stringstyle=\color{codestring},
    commentstyle=\color{codecomment}\itshape,
    numbers=left,
    numberstyle=\tiny\color{gray},
    stepnumber=1,
    numbersep=8pt,
    frame=single,
    rulecolor=\color{codeframe},
    breaklines=true,
    breakatwhitespace=true,
    showstringspaces=false,
    tabsize=4,
    captionpos=b,
    xleftmargin=1.5em,
    framexleftmargin=1.5em
}
\lstdefinestyle{bashstyle}{
    language=bash,
    backgroundcolor=\color{codebg},
    basicstyle=\ttfamily\footnotesize,
    keywordstyle=\color{codekeyword}\bfseries,
    commentstyle=\color{codecomment}\itshape,
    numbers=none,
    frame=single,
    rulecolor=\color{codeframe},
    breaklines=true,
    showstringspaces=false,
    tabsize=2,
    captionpos=b
}
\lstset{style=pythonstyle}

% ---- Diagramas Vectoriales TikZ ----
\usepackage{tikz}
\usetikzlibrary{shapes.geometric, arrows, positioning, calc, fit, backgrounds, patterns}

% ---- Hipervínculos ----
% Limpiar páginas en blanco creadas por \cleardoublepage sin dependencias externas
\makeatletter
\def\cleardoublepage{\clearpage\if@twoside \ifodd\c@page\else
  \hbox{}\thispagestyle{empty}\newpage\if@twocolumn\hbox{}\newpage\fi\fi\fi}
\makeatother

\usepackage{hyperref}
\hypersetup{
    colorlinks=true,
    linkcolor=etsiitblue,
    citecolor=darkgreen,
    urlcolor=etsiitblue,
    filecolor=etsiitblue,
    pdfauthor={Emilio Gullón},
    pdftitle={Análisis predictivo de fatiga mediante datos biométricos multimodales},
    pdfsubject={Trabajo de Fin de Grado - ETSIIT UGR},
    pdfkeywords={Fatiga Mental, Fatiga Física, FatigueSet, Deep Learning, PyTorch, sLSTM, xLSTM, Foundation Models, MOMENT, Chronos, TimesFM, Optuna, Series Temporales}
}

% ---- Listas e Índices ----
\usepackage{enumitem}
\setlist{itemsep=3pt, topsep=4pt}
\usepackage{tocbibind}

% ---- Comandos y Macros Personalizadas ----
\providecommand{\euro}{€}
\newcommand{\fatigueset}{\textsc{FatigueSet}}
\newcommand{\moment}{\textsc{MOMENT}}
\newcommand{\chronos}{\textsc{Chronos}}
\newcommand{\timesfm}{\textsc{TimesFM}}
\newcommand{\xlstm}{\textsc{xLSTM}}
\newcommand{\slstm}{\textsc{sLSTM}}
\newcommand{\mlstm}{\textsc{mLSTM}}
\newcommand{\patchtst}{\textsc{PatchTST}}
\newcommand{\code}[1]{\texttt{#1}}

% ---- Entornos Teorema/Definición Estándar ----
\newtheorem{definicion}{Definición}[chapter]
\newtheorem{hipotesis}{Hipótesis}[chapter]
\newtheorem{observacion}{Observación}[chapter]
